\chapter{Data I/O}
\label{ch:dataio}

Every computation so far has begun and ended inside the session: you typed an
expression, Mathilda evaluated it, and the result scrolled past. This chapter is
about the doorway to the world outside that session---reading numbers a
colleague measured and left in a file, saving a result so the next run can pick
it up, turning a table on disk into a list you can compute with. It is the least
mathematical chapter in the book and, for day-to-day work, one of the most
useful.\index{file I/O}

A word of honesty before we begin. Input and output is a \emph{young} part of
Mathilda, and this chapter is deliberately narrow: it documents what the system
actually does today and is explicit about what it does not. What exists is a
genuine, Mathematica-shaped \emph{stream layer}---you open a file, get a stream
object, and read or write through it---together with \B{ReadList} for reading a
data file field by field, \B{Get} and \B{Put} for saving and reloading whole
Mathilda expressions, and a handful of pure path utilities. What does \emph{not}
yet exist is just as important to know: there is no general \B{Import}/\B{Export}
of data formats (those two are confined to raster images and graphics,
\S\ref{sec:io-importexport}), no binary read/write, and no file-management verbs
such as copying or deleting. Where a gap matters, this chapter names it rather
than papering over it.

Because every example in this book is re-run from scratch at build time (this
chapter is no exception), each one here is self-contained: it writes a
file under \texttt{/tmp} and then reads that same file back, so nothing depends
on data that happens to be lying around. That discipline is worth adopting in
your own scripts, and it makes the round trip---write, then read---the natural
place to start.

\section{The round trip: writing and reading a file}
\label{sec:io-roundtrip}

The smallest complete story in file I/O is to write something and read it back.
Three functions open the door. \B{OpenWrite} opens a file for writing and hands
back a stream object; \B{WriteString} sends raw text to that stream; and \B{Close}
finishes, flushing and releasing the file. To read the data back as numbers,
\mcode{ReadList} does the whole job in one call.

\begin{codepairs}
\pair{09-data-io/roundtrip/1}{\mcode{OpenWrite} truncates the file and returns an
\B{OutputStream} object. The integer is an internal handle---the first stream
opened this session, so \(1\).}
\pair{09-data-io/roundtrip/2}{\mcode{WriteString} sends two lines of text to the
stream; the trailing semicolon suppresses its (empty) result.}
\pair{09-data-io/roundtrip/3}{\mcode{Close} flushes and releases the file,
returning its name to confirm the stream is shut.}
\pair{09-data-io/roundtrip/4}{\B{FileExistsQ} confirms the file is really on
disk.}
\pair{09-data-io/roundtrip/5}{\mcode{ReadList} with the read type \texttt{Number}
scans the whole file and returns its six numbers as a flat list.}
\end{codepairs}

Two things are worth pausing on. First, the newline in the string
\texttt{"88 91 79\textbackslash n95 82 100\textbackslash n"} is a real newline:
Mathilda's string lexer honors the C escapes \mcode{\textbackslash n} and
\mcode{\textbackslash t}, so you can lay out a file's lines directly in a string
literal.\index{string!escape sequences} Second, \mcode{ReadList} did not care that
the six numbers were split across two lines---read as \texttt{Number}, whitespace
and line breaks alike simply separate one token from the next. Where the numbers
sit in the file is a formatting detail; what you get back is the data.

\calloutnote{Under the hood}{A stream reaches the language as an inert object,
\mcode{OutputStream}\texttt{["name", }\emph{id}\texttt{]} or
\B{InputStream}\texttt{["name", }\emph{id}\texttt{]}, whose integer is a handle
into a process-global registry (\texttt{src/io/streams.c}). An output stream
holds an open \texttt{FILE*}; an input stream slurps the whole file into a buffer
with a moving read point. The registry is freed by \mcode{Close} or, for anything
left open, by an \texttt{atexit} hook at program exit, so a session that forgets
to close a stream still leaks nothing under valgrind.}

\usagebox{OpenWrite}

\section{Type-directed reading}
\label{sec:io-readtypes}

The power of \mcode{ReadList} is that you tell it \emph{what kind} of object to read,
and it interprets the bytes accordingly. The same file can be read as numbers, as
words, as whole lines, or as raw bytes---you choose by naming a \emph{read type}.
Take a small table with integers, decimals, and scientific-notation tokens:

\begin{codepairs}
\pair{09-data-io/readtypes/4}{\texttt{Number} reads each token as a number:
integers stay exact, and \texttt{2e5}/\texttt{1.5e-3} become reals.}
\pair{09-data-io/readtypes/5}{\texttt{Real} forces every token to an approximate
number, so even \(1\) comes back as \(1.0\).}
\pair{09-data-io/readtypes/6}{\texttt{Word} keeps each token as a string, doing no
numeric parsing at all.}
\pair{09-data-io/readtypes/7}{\texttt{String} reads a whole line at a time---three
lines, three strings.}
\pair{09-data-io/readtypes/8}{\texttt{Byte} reads raw byte codes; the first five
bytes of \texttt{"1 2 3"} are the ASCII codes \(49, 32, 50, 32, 51\).}
\pair{09-data-io/readtypes/9}{\texttt{Character} reads one-character strings,
byte by byte.}
\pair{09-data-io/readtypes/10}{A third argument caps the read: only the first
three numbers.}
\end{codepairs}

So the same bytes on disk become, depending on the type you ask for, seven
numbers, seven words, three whole lines, or a stream of raw byte codes. The read
type is the lens.\index{ReadList@\mcode{ReadList}!read types}

A \emph{list} of types reads a heterogeneous record: one object of each type per
pass, grouped into a sublist, repeated to end of file. This is exactly how you
read a columnar file whose fields have different types---a name and an age, say:

\begin{codepairs}
\pair{09-data-io/readtypes/14}{\texttt{\{Word, Number\}} reads a string then a
number on each pass, giving one \(\{\text{name}, \text{age}\}\) pair per line.}
\pair{09-data-io/readtypes/18}{When end of file arrives partway through a
pass---here after the lone \(3\)---the unread slot is filled with \B{EndOfFile}
rather than dropped, so the record structure stays intact.}
\end{codepairs}

That \mcode{EndOfFile} in the last pair is not an error; it is how the reader tells
you a record was cut short, so a ragged file never silently loses its shape.

\calloutnote{Theory}{The read types divide cleanly by what counts as a boundary.
\texttt{Byte} and \texttt{Character} consume exactly one byte and never look for a
separator. \texttt{Word}, \texttt{Number}, and \texttt{Real} scan up to the next
\emph{word separator} (whitespace by default). \texttt{Record} and \texttt{String}
scan up to the next \emph{record separator} (a line break). \texttt{Expression}
reads one complete Mathilda expression---the type that makes \mcode{Get} and
\mcode{ReadList} interchangeable for saved results (\S\ref{sec:io-putget}).}

\usagebox{ReadList}

\section{Separators and delimited data}
\label{sec:io-separators}

The word and record separators are not fixed; you set them with options, and that
is what lets \mcode{ReadList} read comma-separated values, tab-separated files, and
other delimited formats. \texttt{WordSeparators} names the strings that break one
token from the next:

\begin{codepairs}
\pair{09-data-io/separators/4}{With \texttt{WordSeparators -> \{","\}}, the comma
becomes the delimiter and a CSV of numbers reads as a flat list.}
\pair{09-data-io/separators/5}{Pair the comma delimiter with a type list to read
the CSV row by row into sublists---a small numeric matrix.}
\end{codepairs}

The complementary option, \texttt{TokenWords}, forces a given string to be split
off as its own word even when nothing surrounds it---useful when an operator or
punctuation mark is glued to its neighbors:

\begin{codepairs}
\pair{09-data-io/separators/9}{Read as plain \texttt{Word}s, \texttt{"a+b"} is a
single token---there is no separator inside it.}
\pair{09-data-io/separators/10}{\texttt{TokenWords -> \{"+"\}} makes the plus sign
a word in its own right, splitting the run into three.}
\end{codepairs}

Alongside these, \texttt{RecordSeparators} sets what ends a line (the default is
the usual trio of \texttt{"\textbackslash r\textbackslash n"},
\texttt{"\textbackslash n"}, and \texttt{"\textbackslash r"}), and the boolean
options \texttt{NullWords} and \texttt{NullRecords} decide whether empty fields
between adjacent separators are kept or skipped---the difference between reading
\texttt{a,,b} as two fields or three.\index{separator!word}\index{separator!record}
Between them, these options cover the common shapes of plain-text tabular
data.

\calloutnote{Pitfall}{Mathilda has no dedicated CSV importer---a comma-delimited
file is just a file you hand the right \texttt{WordSeparators}. That is flexible,
but it also means the quoting and escaping conventions of the CSV standard
(quoted fields containing commas, doubled quotes) are \emph{not} understood: a
field is whatever lies between separators, quotes and all. For clean numeric
tables this is exactly right; for messy real-world CSV with embedded commas, it
is a sharp edge to know about.}

\section{The stream as a cursor}
\label{sec:io-streams}

\mcode{ReadList} reads a whole file in one gulp. When you want to read a file
\emph{incrementally}---a header, then a body sized by what the header said, or a
record at a time in a loop---you open the file yourself and read through the
stream. \B{OpenRead} returns an \mcode{InputStream} that carries a persistent
\emph{current point}, and each \B{Read} advances it, so successive reads return
successive objects.

\begin{codepairs}
\pair{09-data-io/streams/4}{\mcode{OpenRead} returns an \mcode{InputStream}. Its handle
is \(2\): the write stream that created the file was handle \(1\).}
\pair{09-data-io/streams/5}{\mcode{Read} takes one object and moves the point past it.}
\pair{09-data-io/streams/6}{The next \mcode{Read} continues from where the last one
stopped---the cursor persists.}
\pair{09-data-io/streams/7}{A read type may itself be a list: this reads the next
two numbers as a pair.}
\pair{09-data-io/streams/8}{\B{StreamPosition} reports the current point as a byte
offset into the file.}
\pair{09-data-io/streams/9}{\B{Streams} lists every stream still open.}
\pair{09-data-io/streams/10}{\B{SetStreamPosition} moves the point---here back to
the start, returning the new offset.}
\pair{09-data-io/streams/11}{So the very next \mcode{Read} yields the first number
again.}
\end{codepairs}

Positioning turns a stream into random-access storage: \mcode{StreamPosition} tells
you where you are, \mcode{SetStreamPosition} takes you anywhere in the file (and
\texttt{SetStreamPosition[}stream\texttt{, Infinity]} jumps to the end). Because a
type structure can nest, a single \mcode{Read} can pull in a whole shaped object---a
matrix, for instance, read as a list of rows:

\begin{codepairs}
\pair{09-data-io/streams/17}{A nested type structure fills depth-first, so this
reads the file's four numbers into a \(2\times 2\) matrix in one call.}
\pair{09-data-io/streams/21}{Once the point is past the last object, \mcode{Read}
returns \mcode{EndOfFile}---the sentinel a read loop watches for to stop.}
\end{codepairs}

\calloutnote{Under the hood}{\mcode{ReadList} is not a separate engine---it is literally
a loop of the same \texttt{read\_structure()} routine that \mcode{Read} calls once
(\texttt{src/io/read.c}), run to end of file. That is why the two share their read
types and every separator option to the letter: there is one reader, exposed two
ways. The \texttt{Expression} leaf is read \emph{unevaluated} and handed to the
ordinary evaluator, which is why \texttt{Read[s, Expression]} evaluates what it
read while \texttt{Read[s, Hold[Expression]]} keeps it held.}

\section{Persisting expressions: Put, Get, and Write}
\label{sec:io-putget}

The reading so far has treated a file as \emph{data}---numbers and words. But a
file can also hold Mathilda \emph{expressions}, written in the same input syntax
you type at the prompt, and this is how you save a result and reload it in a later
session. \B{Write} sends an expression to a stream in input form; \mcode{Put} (and its
operator \mcode{>>}) writes to a named file; \mcode{Get} reads a file back, evaluating
as it goes.

First, the distinction between the two ways of writing. \mcode{Write} renders each
argument in input form and---crucially---\emph{evaluates} it first, whereas
\mcode{WriteString} writes its text raw:

\begin{codepairs}
\pair{09-data-io/putget/6}{Reading the file back as lines shows what was written:
\mcode{Write}\texttt{[s, 2 + 2]} evaluated to \texttt{4} before writing, \(x^2+1\) was
written in input form, and \mcode{WriteString}'s text went out verbatim.}
\end{codepairs}

For saving results, the \mcode{>>} operator and \mcode{Get} are the everyday tools.
\mcode{>>} is shorthand for \mcode{Put}: it writes one expression to a file, replacing
whatever was there. \mcode{>>>} is \B{PutAppend}: it adds to the file instead.

\begin{codepairs}
\pair{09-data-io/putget/7}{\texttt{FactorInteger[40320] >> file} writes the
factorization of \(8!\) to a file (the result is suppressed).}
\pair{09-data-io/putget/8}{\mcode{Get} reads the file, evaluating each expression, and
returns the value of the last---here the one factorization we wrote.}
\pair{09-data-io/putget/11}{After two \mcode{>>>} appends, \mcode{Get} returns only
the \emph{last} expression in the file: the factorization of \(100\).}
\pair{09-data-io/putget/12}{\mcode{ReadList} reads the same file as
\texttt{Expression}s and \emph{collects} all three, rather than keeping only the
last.}
\end{codepairs}

That contrast between pairs is the whole design. \mcode{Get} is for \emph{loading}: it
runs a file as a script and gives you the final value, exactly as the REPL
bootstrap loads Mathilda's own startup definitions. \mcode{ReadList} is for
\emph{collecting}: it gathers every expression into a list. \mcode{Put} writes several
expressions at once, and they round-trip through \mcode{ReadList} unchanged:

\begin{codepairs}
\pair{09-data-io/putget/13}{\mcode{Put} writes two expressions, one per line.}
\pair{09-data-io/putget/14}{\mcode{ReadList} reads them straight back---a symbolic
result saved to disk and recovered intact.}
\end{codepairs}

\calloutnote{Under the hood}{\mcode{Get} and \mcode{Put} predate the stream layer and live
in \texttt{src/readwrite.c}; \mcode{Get} shares its file-reading core with
\texttt{LoadModule}, the same machinery that loads \texttt{src/internal/init.m} at
startup. Because \mcode{Put} renders with the standard printer, a saved file is
ordinary Mathilda source---you can open it in an editor, and \mcode{Get} it back like
any script.}

\usagebox{Get}

\section{File names and existence}
\label{sec:io-filenames}

The last group of functions is the plumbing around the files themselves: testing
whether a path exists and taking a file name apart. \mcode{FileExistsQ} is the one
that touches the disk; the rest are pure string operations on the path, and never
open anything.

\begin{codepairs}
\pair{09-data-io/filenames/4}{\mcode{FileExistsQ} is \texttt{True} for a file we just
wrote\ldots}
\pair{09-data-io/filenames/5}{\ldots and \texttt{False} for one that is not there.}
\pair{09-data-io/filenames/6}{\B{FileExtension} returns the trailing extension,
without the dot.}
\pair{09-data-io/filenames/7}{\B{FileBaseName} returns the leaf with that one
extension stripped---note it splits off only the last suffix.}
\pair{09-data-io/filenames/10}{\B{FileNameSplit} breaks a path into its
components; the leading \texttt{""} marks an absolute path.}
\pair{09-data-io/filenames/11}{\B{FileNameJoin} is the inverse, assembling
components with the OS separator.}
\pair{09-data-io/filenames/12}{Split then join round-trips to a canonical form of
the original path.}
\end{codepairs}

These are string surgery, portable and side-effect-free---\mcode{FileNameJoin} and
\mcode{FileNameSplit} even take an \texttt{OperatingSystem} option so you can
manipulate Windows paths from Unix or vice versa.\index{file path!components} One
more console convenience rounds out the set: \B{FilePrint} streams the raw bytes
of a file straight to your terminal, the equivalent of the shell's
\texttt{cat}, with optional line ranges (\texttt{FilePrint["f", 1;;10]}). It
writes as a side effect and returns nothing, which is precisely why it does not
appear in the verified transcripts of this chapter---there is no value to
capture. Throughout, we have shown a file's contents by reading them back as data
with \mcode{ReadList}; \mcode{FilePrint} is the quick look when you are at the prompt.

\section{Import and Export: images and graphics}
\label{sec:io-importexport}

A reader coming from Mathematica will expect \mcode{Import} and \mcode{Export} to be the
universal front door for data---CSV, JSON, spreadsheets, and a hundred other
formats. In Mathilda today they are far narrower, and it is important to be plain
about it: \mcode{Import} and \mcode{Export} handle \emph{raster images and graphics
only}. An image round-trips cleanly through a file---

\begin{codepairs}
\pair{09-data-io/importexport/2}{\mcode{Export} writes an image to a PNG and returns
the file name.}
\pair{09-data-io/importexport/3}{\mcode{Import} reads it back; the dimensions survive
the round trip.}
\end{codepairs}

---but a data file is not something they claim, and they say so by declining
rather than guessing:

\begin{codepairs}
\pair{09-data-io/importexport/7}{Handed a plain text file, \mcode{Import} stays
unevaluated---its head comes back as \texttt{Import}, the sign of a call left
alone, not a \mcode{\$Failed}.}
\pair{09-data-io/importexport/8}{Likewise \mcode{Export} declines to write a bare list
to a \texttt{.csv}: there is no data-format writer behind it yet.}
\end{codepairs}

The image and graphics side of \mcode{Import}/\mcode{Export} is real and useful---PNG,
JPEG, BMP, and a resolution-independent vector PDF for plots---and it belongs to
the graphics story, so the graphics chapter (Chapter~\ref{ch:graphics}) develops
it in full. For structured \emph{data}, the honest advice is to reach for the
tools of this chapter instead: \mcode{ReadList} with the right separators for tabular
text, and \mcode{Put}/\mcode{Get} for anything expressible as Mathilda expressions. A
general format layer is future work, not a present feature.

\calloutnote{Pitfall}{The difference between an unevaluated result and
\mcode{\$Failed} is a real signal, not a quirk. \mcode{\$Failed} means ``I tried
and could not''---a missing file, a malformed number. A call that comes back with
its own head, like \texttt{Import[\ldots]} above, means ``I do not handle this at
all,'' leaving the expression untouched so a rule of your own could still act on
it. Reading that distinction saves you from debugging a failure that never
happened.}

\subsection*{Where this connects}

Data I/O is the seam between Mathilda and everything outside it, and its two
halves lead in different directions. The \emph{data} side---\mcode{ReadList} and the
stream layer---feeds the rest of the book: a file of measurements read as a packed
list of numbers is exactly the input the statistics of \S\ref{sec:statistics}
summarize and the array engine of Chapter~\ref{ch:datastructures} accelerates,
and reading a matrix with a nested \mcode{Read} structure produces the same lists that
the linear algebra of \S\ref{sec:linalg} operates on. The \emph{expression}
side---\mcode{Put} and \mcode{Get}---is how a symbolic result outlives its session and how
Mathilda loads its own definitions at startup, a thread the internals chapter
(Chapter~\ref{ch:internals}) picks up. The image and graphics forms of \mcode{Import}
and \mcode{Export} are taken up in Chapter~\ref{ch:graphics}. And for the exhaustive
options of any function named here---every read type, every separator, every
edge case of a path utility---its reference page is one \texttt{?Name} away at the
prompt.
